\chapter{The Fundamental Theorem of Calculus}

The Fundamental Theorem of Calculus shows that two seemingly separate operations are actually inverses of each other.

On one side we have the \emph{derivative}: an operation for measuring
instantaneous rates of change, slopes of tangent lines, and how fast a quantity
changes at a single point. On the other side we have the \emph{integral}: an
operation for accumulating quantities over an interval, measuring
accumulated area, total distance, and swept-out volume. Slope and area, rate and
total, seem like entirely different questions. One concerns local behavior at an instant;
the other concerns global behavior across an interval. One is a tangent line; the other is a region
of the plane.

The Fundamental Theorem of Calculus establishes that these two operations are inverses of
each other. Differentiation and integration undo one another, in the same way that addition
undoes subtraction or squaring undoes square-rooting. The slope question and the
area question are two directions of the same relationship.

That single fact transforms calculus from a collection of clever limit
computations into a systematic method. Before the theorem, finding an area meant
summing infinitely many rectangles and evaluating the limit directly.
After the theorem, finding an area means: find a function whose derivative is
the given integrand, then subtract two of its values. Computations that once
required intricate geometric exhaustion become routine evaluations, and this chapter
explains \emph{why} this connection holds.

\section{What we are bridging}

Let us set the two sides down precisely so we can see how far apart they start.

\subsection{The integral, briefly recalled}

Given a function $f$ on an interval $[a,b]$, the \emph{definite integral}
\[
  \int_a^b f(x)\,\dd x
\]
is the limit of Riemann sums. You chop $[a,b]$ into $n$ little pieces, build a
thin rectangle on each piece whose height is a value of $f$ there, add up the
rectangle areas, and let the pieces shrink to nothing:
\[
  \int_a^b f(x)\,\dd x \;=\; \lim_{n\To\infty}\sum_{k=1}^{n} f(x_k^*)\,\Delta x_k .
\]
When $f\ge 0$ this is the area under the graph. When $f$ dips below the axis it
is signed area -- area above the axis counts positive, area below counts
negative. The key point for us: this is a \emph{global} object. To know it you
must, in principle, account for the value of $f$ at \emph{every} point of
$[a,b]$.

\subsection{The derivative, briefly recalled}

Given a function $F$, its \emph{derivative} at a point $x$ is the local rate of
change,
\[
  F'(x) \;=\; \lim_{h\To 0}\frac{F(x+h)-F(x)}{h},
\]
the slope of the tangent line to the graph of $F$ at $x$. This is a
\emph{local} object: it cares only about the behavior of $F$ in an
infinitesimally small neighborhood of the one point $x$.

\begin{intuitionbox}
Here is the core question. The integral depends on values across an entire
interval. The derivative depends only on behavior in the immediate neighborhood of
a single point. How can an interval-wide accumulation (accumulated area) and
an instantaneous rate of change (slope) determine each other?

The resolution is the central image of this chapter: \textbf{watch the area
grow}. Don't compute the area once and for all. Instead, let the right-hand edge
slide, and ask how fast the accumulated area is increasing at each moment. That
\emph{rate} of accumulation is a derivative -- and it turns out to equal the
height of the curve. Accumulating and rate-measuring are two directions of a
single relationship.
\end{intuitionbox}

\section{The area-so-far function}

Everything hinges on a single new idea: instead of treating the area as a
\emph{number}, treat it as a \emph{function}. Fix the left endpoint $a$ and let
the right endpoint vary. Define the \emph{area-so-far} function (the standard
name is the \emph{accumulation function}):
\[
  A(x) \;=\; \int_a^x f(t)\,\dd t .
\]

\begin{focusbox}
Read $A(x)$ as ``the signed area under $f$ from the starting line $a$ out to the
moving wall at $x$.'' As you push the wall $x$ to the right, you sweep out more
area, so $A$ grows. As you pull the wall back left, $A$ shrinks. $A(x)$ is a
running total.

Notice the bound of integration is now $x$, and the variable inside the integral
had to be renamed $t$ so it doesn't collide with the limit. The $t$ is a dummy --
it gets summed over and disappears. The surviving variable, the one $A$ actually
depends on, is the upper limit $x$.
\end{focusbox}

\begin{figure}[H]
\centering
\begin{tikzpicture}
\begin{axis}[
  width=11cm, height=6.4cm,
  axis lines=left,
  xlabel={$t$}, ylabel={$f(t)$},
  xtick={1,3}, xticklabels={$a$,$x$},
  ytick=\empty,
  xmin=0, xmax=5, ymin=0, ymax=3.2,
  domain=0:5, samples=120,
  clip=true,
]
\addplot[draw=none, fill=areafill, opacity=0.6, domain=1:3] {0.6+0.35*x+0.08*x^2} \closedcycle;
\addplot[curveblue, very thick, domain=0:5] {0.6+0.35*x+0.08*x^2};
\node at (axis cs:2,1.0) {$A(x)$};
\draw[dashed] (axis cs:3,0) -- (axis cs:3,2.37);
\draw[dashed] (axis cs:1,0) -- (axis cs:1,1.03);
\end{axis}
\end{tikzpicture}
\caption{The accumulation function $A(x)=\int_a^x f(t)\,\dd t$ is the shaded area
from the fixed wall $a$ to the sliding wall $x$. Sliding $x$ to the right grows
the area.}
\label{ch21:fig-accum}
\end{figure}

\begin{labbox}
Take the simplest non-trivial example: $f(t)=t$ on $[0,x]$. The region under
$f$ from $0$ to $x$ is a right triangle with both legs of length $x$, so its
area is $\tfrac12 x^2$. Hence
\[
  A(x) = \int_0^x t\,\dd t = \tfrac12 x^2 .
\]
Now do the central calculation of this chapter: differentiate the area.
\[
  A'(x) = \ddx{\tfrac12 x^2} = x = f(x).
\]
The derivative of the area-so-far is the original height. That is not an isolated
property of triangles: it is the Fundamental Theorem in its simplest setting.
Hold onto this result; we will explain it in full next.
\end{labbox}

\begin{recoverybox}
\textbf{Pause and think:} If $f(t) > 0$ everywhere on $[a,b]$, can the
accumulation function $A(x)=\int_a^x f(t)\,\dd t$ ever decrease as $x$ increases?
Picture the sliding wall: whenever you advance $x$ to the right, the newly swept
strip sits above the horizontal axis, adding strictly positive area. So $A(x)$
must be strictly increasing. An increasing function has a non-negative
derivative: indeed, $A'(x) = f(x) > 0$. The sign of the curve directly controls
the slope of the running area.
\end{recoverybox}

\section{Part 1: differentiating the area gives back the curve}

\subsection{The infinitesimal-strip picture}

Before any $\eps$ appears, let us see \emph{why} $A'(x)=f(x)$ has to be true,
using a simple geometric argument.

Push the right wall from $x$ to $x+h$, where $h$ is a tiny positive number. The
area grows by a sliver: the strip sitting between $x$ and $x+h$. Call that extra
area $A(x+h)-A(x)$. What does the sliver look like? It is a thin vertical strip
of width $h$, sitting under the curve $f$.

\begin{figure}[H]
\centering
\begin{tikzpicture}
% Macro view (left)
\begin{axis}[
  name=macro,
  width=7.2cm, height=5.5cm,
  axis lines=left,
  xlabel={$t$}, ylabel={$f(t)$},
  title style={font=\small},
  title={\textbf{(a)} Accumulation and slice of width $h$},
  xtick={1.0,2.4,3.1}, xticklabels={$a$,$x$,$x{+}h$},
  ytick=\empty,
  xmin=0.5, xmax=4.2, ymin=0, ymax=4.2,
  domain=0.5:4.0, samples=100,
  clip=false,
]
\addplot[draw=none, fill=areafill, opacity=0.35, domain=1.0:2.4] {0.7+0.3*x+0.22*x^2} \closedcycle;
\addplot[draw=none, fill=quotientorange, opacity=0.55, domain=2.4:3.1] {0.7+0.3*x+0.22*x^2} \closedcycle;
\addplot[curveblue, very thick, domain=0.8:3.8] {0.7+0.3*x+0.22*x^2};
\draw[dashed, gray] (axis cs:2.4,0) -- (axis cs:2.4,2.69);
\draw[dashed, gray] (axis cs:3.1,0) -- (axis cs:3.1,3.74);
\node at (axis cs:1.7,1.0) {$A(x)$};
\node[quotientorange!80!black, font=\footnotesize] at (axis cs:2.75,1.5) {$\Delta A$};
\end{axis}

% Micro zoom view (right)
\begin{axis}[
  name=micro,
  at={($(macro.east)+(1.8cm,0)$)},
  anchor=west,
  width=6.8cm, height=5.5cm,
  axis lines=left,
  xlabel={}, ylabel={},
  title style={font=\small},
  title={\textbf{(b)} Zoom: trapped by min and max},
  xtick={2.4,3.1}, xticklabels={$x$,$x{+}h$},
  ytick={2.69,3.74}, yticklabels={$m_h$,$M_h$},
  xmin=2.2, xmax=3.3, ymin=0, ymax=4.2,
  domain=2.4:3.1, samples=60,
  clip=false,
]
% Inscribed rectangle
\addplot[draw=none, fill=targetgreen, opacity=0.25, domain=2.4:3.1] {2.69} \closedcycle;
% Extra strip area under curve
\addplot[draw=none, fill=quotientorange, opacity=0.45, domain=2.4:3.1] {0.7+0.3*x+0.22*x^2} \closedcycle;
% Circumscribed rectangle top line
\draw[dashed, red!70!black, thick] (axis cs:2.4,3.74) -- (axis cs:3.1,3.74);
% Curve
\addplot[curveblue, very thick, domain=2.3:3.2] {0.7+0.3*x+0.22*x^2};
\draw[dashed, gray] (axis cs:2.4,0) -- (axis cs:2.4,2.69);
\draw[dashed, gray] (axis cs:3.1,0) -- (axis cs:3.1,3.74);
% Annotations
\draw[<->, thick] (axis cs:2.4,-0.55) -- (axis cs:3.1,-0.55) node[midway, below, font=\small]{$h$};
\node[targetgreen!60!black, font=\footnotesize, anchor=west] at (axis cs:2.45,1.2) {$m_h \cdot h \le \Delta A$};
\node[red!70!black, font=\footnotesize, anchor=south east] at (axis cs:3.1,3.8) {$\le M_h \cdot h$};
\end{axis}

% Zoom dashed connector lines between lanes
\draw[dashed, gray!60, thick] (macro.south -| macro.east) -- (micro.south -| micro.west);
\draw[dashed, gray!60, thick] (macro.north -| macro.east) -- (micro.north -| micro.west);

\end{tikzpicture}
\caption{The rate of accumulation. \textbf{(a)} Growing the wall from $x$ to $x+h$ adds the thin sliver $\Delta A = A(x+h)-A(x)$. \textbf{(b)} Zooming in on the sliver reveals it trapped between the lower rectangle of area $m_h \cdot h$ and the upper rectangle of area $M_h \cdot h$. As $h \to 0$, both bounding heights converge to $f(x)$, squeezing the average rate $\Delta A / h$ to $f(x)$.}
\label{ch21:fig-strip}
\end{figure}

Here is the key geometric observation: if $h$ is small, the curve barely changes
across the strip, so the strip is \emph{almost a rectangle}: width $h$, height
roughly $f(x)$. Its area is approximately
\[
  A(x+h)-A(x) \;\approx\; f(x)\cdot h .
\]
Divide both sides by $h$:
\[
  \frac{A(x+h)-A(x)}{h} \;\approx\; f(x).
\]
The left side is exactly the difference quotient whose limit defines $A'(x)$. As
$h\To 0$ the strip gets thinner, the ``almost a rectangle'' becomes ``a
rectangle'' in the limit, the approximation tightens to equality, and we land on
\[
  A'(x) = f(x).
\]

\begin{happeningbox}
What is really happening when you write $A'(x)=f(x)$? You are reading off the
\textbf{instantaneous rate at which area is piling up.} The height of the curve
at $x$ \emph{is} the speed at which the running-area total grows as the wall
sweeps past $x$.

Think of $f$ as a flow rate -- liters per second pouring into a tank -- and
$A(x)$ as the total liters in the tank after time $x$. ``How fast is the tank
filling right now?'' is answered by the current flow rate. That sentence is the
Fundamental Theorem, Part 1, in plain language: \emph{the rate of accumulation is
the integrand itself}. A thin strip represents the contribution across a narrow increment;
its height is the flow rate, its width $\dd x$ is the increment of time, and the product
$f(x)\,\dd x$ is the incremental area added across that step.
\end{happeningbox}

\subsection{Making it a real proof}

The strip picture is convincing, but ``$\approx$'' is a promise, not a proof. The
real worry is this: across the strip the height is not exactly $f(x)$; the
curve wobbles. We need to control that wobble and show it vanishes in the limit.
The clean tool for that is the average-value squeeze, and it works for any
continuous $f$.

\begin{theorem}[Fundamental Theorem of Calculus, Part 1]\label{ch21:thm-ftc1}
Let $f$ be continuous on $[a,b]$. Define $A:[a,b]\To\R$ by
\[
  A(x) = \int_a^x f(t)\,\dd t .
\]
Then $A$ is differentiable on $(a,b)$ (and one-sidedly at the endpoints), and
\[
  A'(x) = f(x) \qquad\text{for every } x .
\]
In words: the accumulation function is an antiderivative of $f$.
\end{theorem}

\begin{intuitionbox}
The idea of the proof is to trap the increment $A(x+h)-A(x)$ -- the area of the
strip -- between the area of the shortest rectangle that fits under it and the
area of the tallest rectangle that covers it. Both rectangles have width $h$;
their heights are the minimum and maximum of $f$ on the little interval. As
$h\To 0$ that little interval collapses onto the point $x$, so both the minimum
and the maximum of $f$ are squeezed toward $f(x)$ (this is where continuity earns
its keep). The strip's average height has nowhere to go but $f(x)$.
\end{intuitionbox}

\begin{proof}
Fix $x$ and take $h>0$ small enough that $x+h$ stays in $[a,b]$. Using the
splitting rule for integrals (area from $a$ to $x+h$ equals area from $a$ to $x$
plus area from $x$ to $x+h$),
\[
  A(x+h)-A(x)
  = \int_a^{x+h} f - \int_a^{x} f
  = \int_x^{x+h} f(t)\,\dd t .
\]
So the increment of $A$ is exactly the area of the strip -- no approximation yet,
this is an identity.

Because $f$ is continuous on the closed interval $[x,x+h]$, the Extreme Value
Theorem (which guarantees that any continuous curve on a closed interval
actually reaches its lowest trough and highest peak) hands us a point where $f$
attains its minimum value $m_h$ and a point where it attains its maximum value
$M_h$ on that interval. Every value of $f$ on $[x,x+h]$ lies between $m_h$ and
$M_h$, so the strip's area is trapped between the two rectangles of width $h$:
\[
  m_h\, h \;\le\; \int_x^{x+h} f(t)\,\dd t \;\le\; M_h\, h .
\]
Divide through by $h>0$:
\[
  m_h \;\le\; \frac{A(x+h)-A(x)}{h} \;\le\; M_h .
\]
Now let $h\To 0^+$. The interval $[x,x+h]$ shrinks to the single point $x$. By
continuity of $f$, both the minimum value $m_h$ and the maximum value $M_h$
converge to $f(x)$: any point achieving them is within $h$ of $x$, and
$f(\text{point})\To f(x)$. The middle term is squeezed between two quantities
both heading to $f(x)$, so by the Squeeze Theorem
\[
  \lim_{h\To 0^+}\frac{A(x+h)-A(x)}{h} = f(x).
\]
The case $h<0$ follows the exact same logic with signed directions kept straight.
When $h<0$, the interval of integration is $[x+h,x]$ of positive length
$-h=\abs{h}$. The change in area is $A(x+h)-A(x) = -\int_{x+h}^x f(t)\,\dd t$.
Bounding $f$ between its minimum $m_h$ and maximum $M_h$ on $[x+h,x]$ gives
$m_h(-h) \le \int_{x+h}^x f(t)\,\dd t \le M_h(-h)$. Negating and dividing by
the negative quantity $h$ reverses the inequalities right back to
$m_h \le \frac{A(x+h)-A(x)}{h} \le M_h$. Letting $h\To 0^-$ squeezes the
left-hand limit to $f(x)$ as well. The two one-sided limits agree, so the true
two-sided derivative $A'(x)$ exists and equals $f(x)$.
\end{proof}

\begin{whybox}
Why does the proof go through the \emph{min} and \emph{max} on the little
interval rather than just using $f(x)$? Because we are not allowed to assume the
strip is a perfect rectangle -- that is the very thing we are trying to justify.
The min and max rectangles are airtight bounds we can write down with no
hand-waving: the true strip area genuinely sits between them. Continuity then
guarantees the trap is tight, because over a shrinking interval the highest and
lowest values of a continuous function cannot stay apart -- they both must
collapse to the value at the point. Continuity is exactly the hypothesis that
makes ``the curve barely changes across a thin strip'' a theorem instead of a
hope.
\end{whybox}

\begin{remark}
There is an equivalent and very clean way to package the squeeze, called the
\emph{Mean Value Theorem for Integrals}: for continuous $f$ there is a point
$c\in[x,x+h]$ with $\int_x^{x+h} f = f(c)\cdot h$. (The average height of the
strip is actually achieved somewhere by $f$.) Then
$\frac{A(x+h)-A(x)}{h}=f(c)$, and as $h\To 0$ the point $c$ is trapped at $x$, so
$f(c)\To f(x)$. This expresses the same squeeze argument in concise form.
\end{remark}

\begin{trapbox}
A frequent confusion: students see $\int_a^x f(t)\,\dd t$ and try to ``plug in''
or differentiate with respect to $t$. There is no $t$ to differentiate with
respect to -- $t$ has been integrated away. The function $A$ depends on its
\emph{upper limit} $x$. The statement $A'(x)=f(x)$ means: differentiate with
respect to the upper limit, and you recover the integrand evaluated there.

Also resist writing $A(x)=\int_a^x f(x)\,\dd x$ with $x$ in both places. Using
the same letter for the dummy variable and the limit is a recipe for nonsense.
Always rename the inside variable.
\end{trapbox}

\subsection{The chain-rule version (variable in the limits)}

Part 1 plus the chain rule lets you differentiate any accumulation whose limits
are themselves functions. Suppose
\[
  G(x) = \int_a^{u(x)} f(t)\,\dd t .
\]
Write $G(x)=A(u(x))$ where $A$ is the accumulation function from before, so
$A'=f$. The chain rule gives
\[
  G'(x) = A'(u(x))\cdot u'(x) = f\big(u(x)\big)\, u'(x).
\]
If the variable sits in the lower limit instead, flip the sign (swapping limits
negates the integral). The fully general rule, with both limits moving:
\[
  \ddx{\int_{u(x)}^{v(x)} f(t)\,\dd t}
  = f\big(v(x)\big)\,v'(x) - f\big(u(x)\big)\,u'(x).
\]

\begin{example}\label{ch21:ex-chain}
Differentiate $\displaystyle G(x)=\int_2^{x^3}\sin(t^2)\,\dd t$.

\textbf{Solution.} Here $f(t)=\sin(t^2)$ and the upper limit is $u(x)=x^3$, so
$u'(x)=3x^2$. By the chain-rule version,
\[
  G'(x) = \sin\big((x^3)^2\big)\cdot 3x^2 = 3x^2\sin(x^6).
\]
Notice we never found a formula for $G$ itself -- there isn't an elementary one,
since $\sin(t^2)$ has no elementary antiderivative. Yet we differentiated $G$
exactly. This is what Part 1 buys you: the derivative of an integral, with no
need to evaluate the integral first.
\end{example}

\begin{example}\label{ch21:ex-chain2}
Differentiate $\displaystyle H(x)=\int_{x}^{x^2}\frac{1}{1+\ln t}\,\dd t$ for
$x>1$.

\textbf{Solution.} Both limits move: $u(x)=x$, $v(x)=x^2$. With
$f(t)=\frac{1}{1+\ln t}$,
\[
  H'(x) = \frac{1}{1+\ln(x^2)}\cdot 2x \;-\; \frac{1}{1+\ln x}\cdot 1
        = \frac{2x}{1+2\ln x} - \frac{1}{1+\ln x}.
\]
\end{example}

\section{Part 1 says: every continuous function has an antiderivative}

Tucked inside Theorem~\ref{ch21:thm-ftc1} is a statement of pure existence that
is easy to walk past but worth pausing on.

\begin{corollary}[Existence of antiderivatives]\label{ch21:cor-exist}
Every function continuous on an interval has an antiderivative on that interval --
namely its own accumulation function $A(x)=\int_a^x f(t)\,\dd t$.
\end{corollary}

\begin{aestheticbox}
This existence statement is stronger than it might appear. In introductory calculus, ``find an antiderivative of $f$'' often
means ``find a familiar formula whose derivative is $f$,'' and for many ordinary
functions -- $e^{-x^2}$, $\frac{\sin x}{x}$, $\sqrt{1+x^3}$ -- \emph{no such
formula exists}. The absence of an elementary formula is not a limitation of technique;
Liouville proved that these antiderivatives simply cannot be expressed in terms of
elementary functions.

Part 1 establishes that the antiderivative still \emph{exists} as a well-defined function; we define it by its accumulation
integral rather than by an algebraic formula. The Gaussian integral,
$\int_0^x e^{-t^2}\,\dd t$, is essential in probability, where (up to a factor of $2/\sqrt{\pi}$) it defines the error function $\operatorname{erf}$. It is a fully legitimate function with a
well-defined derivative ($e^{-x^2}$, by Part~1).
\end{aestheticbox}

\section{Part 2: evaluating an integral with any antiderivative}

Part 1 tells us \emph{a} particular antiderivative -- the area function $A$ --
satisfies $A'=f$. Part 2 turns this into an efficient computational method: to find the
area, you may use \emph{any} antiderivative you can find, and the
answer reduces to a simple subtraction.

\begin{theorem}[Fundamental Theorem of Calculus, Part 2]\label{ch21:thm-ftc2}
Let $f$ be continuous on $[a,b]$ and let $F$ be \emph{any} antiderivative of $f$
on $[a,b]$, meaning $F'=f$ throughout. Then
\[
  \int_a^b f(x)\,\dd x \;=\; F(b)-F(a).
\]
\end{theorem}

The standard shorthand for the right side is $\eval{F(x)}{a}^{b}=F(b)-F(a)$, read
``$F$ evaluated from $a$ to $b$.''

\begin{intuitionbox}
Why should the area equal $F(b)-F(a)$ for \emph{any} antiderivative, not just the
area function? Because antiderivatives of the same function can differ only by a
constant (by the Mean Value Theorem, two functions with
the same derivative maintain a constant vertical gap everywhere). The area function $A$
and your chosen $F$ differ by a fixed constant $C$. When you compute the
difference $F(b)-F(a)$, the constant cancels: $(A(b)+C)-(A(a)+C)=A(b)-A(a)$. And
$A(b)-A(a)$ is the total area from $a$ to $b$. The constant $+C$ cancels
completely in the subtraction.
\end{intuitionbox}

\begin{proof}[Proof via Part 1]
Let $A(x)=\int_a^x f(t)\,\dd t$ be the area function. By Part 1, $A'=f$. By
assumption $F'=f$ as well. Two functions with the same derivative on an interval
differ by a constant (vanishing-derivative corollary of the MVT), so there is a
constant $C$ with
\[
  F(x) = A(x) + C \qquad\text{for all } x\in[a,b].
\]
Now simply evaluate the difference:
\[
  F(b)-F(a) = \big(A(b)+C\big)-\big(A(a)+C\big) = A(b)-A(a).
\]
But $A(a)=\int_a^a f = 0$ (an interval of zero width sweeps out no area) and
$A(b)=\int_a^b f$. Therefore $F(b)-F(a)=\int_a^b f(x)\,\dd x$.
\end{proof}

There is a second proof that does not lean on Part 1 at all -- it grows the total
out of tiny MVT-certified increments. It is worth seeing, because it makes the
telescoping cancellation that powers Part 2 completely visible.

\begin{proof}[Proof via telescoping and the Mean Value Theorem]
Partition $[a,b]$ into pieces $a=x_0<x_1<\dots<x_n=b$. The total change in $F$
across $[a,b]$ telescopes -- each interior value is added once and subtracted
once, so all that survives is the endpoints:
\[
  F(b)-F(a) = \sum_{k=1}^{n}\big(F(x_k)-F(x_{k-1})\big).
\]
Apply the Mean Value Theorem to $F$ on each little interval $[x_{k-1},x_k]$:
there is a sample point $c_k$ inside with
\[
  F(x_k)-F(x_{k-1}) = F'(c_k)\,(x_k-x_{k-1}) = f(c_k)\,\Delta x_k .
\]
(We used $F'=f$.) Substituting,
\[
  F(b)-F(a) = \sum_{k=1}^{n} f(c_k)\,\Delta x_k .
\]
The right side is a genuine Riemann sum for $f$ -- with the MVT choosing the
sample points for us. As the partition is refined and the mesh goes to $0$, the
left side does not move at all (it is just $F(b)-F(a)$, a fixed number), while the
right side converges to $\int_a^b f$ by definition of the integral (which exists
and is partition-independent because $f$ is continuous). Therefore
$F(b)-F(a)=\int_a^b f(x)\,\dd x$.
\end{proof}

\begin{whybox}
The second proof exposes the true mechanism of Part 2: \textbf{a total change is
a sum of little changes, and a little change in $F$ is, by the MVT, exactly a
height of $f$ times a width.} Stacking those little $f(c_k)\,\Delta x_k$ bricks
\emph{is} building the area. Telescoping is what guarantees those bricks reassemble
into the clean endpoint difference $F(b)-F(a)$ -- every interior wall is shared by
two bricks and cancels, leaving only the outer walls $a$ and $b$ (Figure~\ref{ch21:fig-telescoping}). Integration is
addition of infinitesimal increments; an antiderivative is a running total of
them; so the net total is the final reading minus the initial reading. Subtraction
of two antiderivative values is just the bookkeeping of ``ending balance minus
starting balance.''
\end{whybox}

\begin{figure}[H]
\centering
\begin{tikzpicture}
% Top lane: integrand and Riemann slices
\begin{axis}[
  name=toplane,
  width=11.0cm, height=4.6cm,
  axis lines=left,
  xlabel={$x$}, ylabel={$f(x)$},
  xlabel style={at={(ticklabel* cs:1.01)}, anchor=west},
  title style={font=\small},
  title={\textbf{Top Lane:} Slicing the rate curve $f(x)$ into area increments $f(c_k)\Delta x_k$},
  xtick={1, 2, 3, 4, 5},
  xticklabels={$a{=}x_0$, $x_1$, $x_2$, $x_3$, $b{=}x_4$},
  ytick=\empty,
  xmin=0.5, xmax=5.5, ymin=0, ymax=4.2,
  domain=0.8:5.2, samples=80,
  clip=false,
]
% 4 sample rectangles
\addplot[draw=sourceblue!80!black, fill=sourceblue!25, thick] coordinates {(1,0) (1,1.5) (2,1.5) (2,0)} -- cycle;
\addplot[draw=targetgreen!80!black, fill=targetgreen!25, thick] coordinates {(2,0) (2,2.2) (3,2.2) (3,0)} -- cycle;
\addplot[draw=quotientorange!80!black, fill=quotientorange!25, thick] coordinates {(3,0) (3,3.1) (4,3.1) (4,0)} -- cycle;
\addplot[draw=structurepurple!80!black, fill=structurepurple!25, thick] coordinates {(4,0) (4,3.8) (5,3.8) (5,0)} -- cycle;

% Continuous curve f(x)
\addplot[curveblue, very thick, domain=0.8:5.2] {0.8 + 0.6*x};

\node[font=\footnotesize, sourceblue!80!black] at (axis cs:1.5,0.75) {$\Delta F_1$};
\node[font=\footnotesize, targetgreen!80!black] at (axis cs:2.5,1.1) {$\Delta F_2$};
\node[font=\footnotesize, quotientorange!80!black] at (axis cs:3.5,1.55) {$\Delta F_3$};
\node[font=\footnotesize, structurepurple!80!black] at (axis cs:4.5,1.9) {$\Delta F_4$};
\end{axis}

% Bottom lane: antiderivative F(x) rising step by step
\begin{axis}[
  name=botlane,
  at={($(toplane.south)-(0,2.1cm)$)},
  anchor=north,
  width=11.0cm, height=4.6cm,
  axis lines=left,
  xlabel={$x$}, ylabel={$F(x)$},
  title style={font=\small, yshift=2pt},
  title={\textbf{Bottom Lane:} Stacking vertical changes: $\sum \Delta F_k = F(b)-F(a)$ telescopes},
  xtick={1, 2, 3, 4, 5},
  xticklabels={$a$, $x_1$, $x_2$, $x_3$, $b$},
  ytick={0.8, 11.4},
  yticklabels={$F(a)$, $F(b)$},
  xmin=0.5, xmax=5.5, ymin=0, ymax=12.5,
  domain=0.8:5.2, samples=80,
  clip=false,
]
% Staircase rises
\draw[very thick, sourceblue!80!black] (axis cs:1,0.8) -- (axis cs:2,0.8) -- (axis cs:2,2.3);
\draw[very thick, targetgreen!80!black] (axis cs:2,2.3) -- (axis cs:3,2.3) -- (axis cs:3,4.5);
\draw[very thick, quotientorange!80!black] (axis cs:3,4.5) -- (axis cs:4,4.5) -- (axis cs:4,7.6);
\draw[very thick, structurepurple!80!black] (axis cs:4,7.6) -- (axis cs:5,7.6) -- (axis cs:5,11.4);

% Smooth F(x) curve
\addplot[tangentred, very thick, dashed, domain=1:5] {0.8*x + 0.3*x^2 - 0.3};

% Overall rise bracket
\draw[decorate, decoration={brace, amplitude=6pt, mirror}, thick] (axis cs:5.15,0.8) -- (axis cs:5.15,11.4)
  node[midway, right=8pt, font=\small] {Total rise: $F(b)-F(a)$};

\draw[dashed, gray] (axis cs:1,0.8) -- (axis cs:5.15,0.8);
\draw[dashed, gray] (axis cs:5,11.4) -- (axis cs:5.15,11.4);
\end{axis}

% Connector arrows in whitespace between lanes
\draw[-{Stealth[length=2.2mm]}, sourceblue, thick] ([yshift=-0.4cm]$(toplane.south west)!0.22!(toplane.south east)$) -- ([yshift=1.0cm]$(botlane.north west)!0.22!(botlane.north east)$);
\draw[-{Stealth[length=2.2mm]}, targetgreen, thick] ([yshift=-0.4cm]$(toplane.south west)!0.42!(toplane.south east)$) -- ([yshift=1.0cm]$(botlane.north west)!0.42!(botlane.north east)$);
\draw[-{Stealth[length=2.2mm]}, quotientorange, thick] ([yshift=-0.4cm]$(toplane.south west)!0.62!(toplane.south east)$) -- ([yshift=1.0cm]$(botlane.north west)!0.62!(botlane.north east)$);
\draw[-{Stealth[length=2.2mm]}, structurepurple, thick] ([yshift=-0.4cm]$(toplane.south west)!0.82!(toplane.south east)$) -- ([yshift=1.0cm]$(botlane.north west)!0.82!(botlane.north east)$);

\end{tikzpicture}
\caption{The telescoping mechanism of FTC Part 2. \textbf{Top lane:} Slices of area under the rate curve $f(x)$ have sizes $f(c_k)\Delta x_k$. \textbf{Bottom lane:} By the Mean Value Theorem, each area slice equals an exact vertical rise $\Delta F_k = F(x_k)-F(x_{k-1})$ on the antiderivative curve. Summing the slices stacks the rises, and internal heights cancel, leaving only the net rise $F(b)-F(a)$.}
\label{ch21:fig-telescoping}
\end{figure}

\begin{happeningbox}
The cleanest way to feel Part 2 is the bank-account reading. Let $F(t)$ be your
account balance at time $t$ and $f(t)=F'(t)$ the instantaneous rate money flows
in (deposits positive, withdrawals negative). Then $\int_a^b f(t)\,\dd t$ adds up
all the little flows over the period -- the net amount that moved -- and that must
equal the change in balance, $F(b)-F(a)$. You did not need to track every
transaction; you only needed the balance at the start and the balance at the end.
The Fundamental Theorem is the statement that \emph{net accumulated rate equals
net change in total}, and it is so intuitive in this form that it feels like it
shouldn't need a proof. The work of the chapter is showing that the intuitive
banking fact and the geometric area fact are the very same theorem.
\end{happeningbox}

\section{Why differentiation and integration are inverse}

We can now say cleanly what ``inverse operations'' means here. There are two
round trips, one for each order.

\begin{focusbox}
\textbf{Integrate, then differentiate} (Part 1): start with $f$, accumulate it
into $A(x)=\int_a^x f$, then differentiate. You get $f$ back.
\[
  \ddx{\int_a^x f(t)\,\dd t} = f(x).
\]

\textbf{Differentiate, then integrate} (Part 2): start with $F$, differentiate
into $F'=f$, then accumulate. You recover the net change in $F$.
\[
  \int_a^b F'(x)\,\dd x = F(b)-F(a).
\]

Each operation undoes the other -- with the single wrinkle that going
$F\To F'\To \int F'$ recovers $F$ only \emph{up to its value at the starting
point} (you get the change $F(b)-F(a)$, not $F$ on the nose). That missing piece
is the constant of integration, and it is the only thing the two operations
cannot quite cancel.
\end{focusbox}

\begin{figure}[H]
\centering
\begin{tikzpicture}[
  box/.style={draw=structurepurple, very thick, rounded corners, fill=structurepurple!8,
              minimum width=2.6cm, minimum height=1.1cm, font=\large},
]
\node[box] (f) at (0,0) {$f$};
\node[box] (F) at (6,0) {$F,\ F'=f$};
\draw[-{Stealth[length=3mm]}, very thick, curveblue]
  (f.north east) to[bend left=22]
  node[above=4pt, font=\small, text=curveblue]{integrate $\int_a^x$} (F.north west);
\draw[-{Stealth[length=3mm]}, very thick, tangentred]
  (F.south west) to[bend left=22]
  node[below, font=\small, text=tangentred]{differentiate $\dv{}{x}$} (f.south east);
\end{tikzpicture}
\caption{The two operations form a loop. Integrating $f$ produces an
antiderivative $F$; differentiating $F$ returns $f$. The Fundamental Theorem says
the loop closes.}
\label{ch21:fig-loop}
\end{figure}

\begin{connectionbox}
This is the same pattern you meet whenever an operation has an inverse, but here
it operates on whole \emph{functions} rather than numbers. Differentiation $D$
and integration are like a matrix and its inverse, or like a group element and
its inverse: $D$ followed by $\int$ is (almost) the identity, and $\int$ followed
by $D$ is the identity. The ``almost'' -- the loss of the constant -- is exactly
why $D$ has a one-dimensional kernel (in linear algebra, the kernel is the subspace
of inputs crushed to zero): here, the constant functions, whose derivative is
zero, are the information differentiation throws away and integration cannot
restore. You will meet this kernel-and-inverse bookkeeping again in linear
algebra and again in the theory of differential equations, where the $+C$
generalizes to the full family of solutions.
\end{connectionbox}

\section{Antiderivatives and the $+C$}

An \emph{antiderivative} of $f$ is any function $F$ with $F'=f$. The
\emph{indefinite integral} collects them all:
\[
  \int f(x)\,\dd x = F(x) + C ,
\]
where $F$ is one antiderivative and $C$ ranges over all real constants. The
notation deliberately drops the limits to signal ``a family of functions,'' as
opposed to the definite integral $\int_a^b f$, which is ``a single number.''

\begin{whybox}
Why is the family \emph{exactly} $F(x)+C$ -- why no other antiderivatives lurking
elsewhere? Suppose $G'=f$ too. Then $(G-F)'=f-f=0$ on the interval. A function
with zero derivative on an interval is constant (Mean Value Theorem again: if it
ever changed, some secant slope would be nonzero, forcing a nonzero derivative
somewhere). So $G-F=C$, i.e. $G=F+C$. The constant $+C$ gives the
\emph{complete} description of all antiderivatives on an interval. Geometrically: all
antiderivatives are vertical shifts of one another -- same shape, same slope at
every $x$, just slid up or down. Shifting a graph vertically never changes any
slope, so it never changes the derivative.
\end{whybox}

\begin{figure}[H]
\centering
\begin{tikzpicture}
\begin{axis}[
  width=10.5cm, height=6cm,
  axis lines=middle,
  xlabel={$x$}, ylabel={$y$},
  xmin=-2.4, xmax=2.4, ymin=-2.2, ymax=4.4,
  domain=-2.1:2.1, samples=80,
  xtick=\empty, ytick=\empty, clip=false,
]
\addplot[curveblue, very thick] {x^2/2 - 1};
\addplot[curveblue, very thick] {x^2/2};
\addplot[curveblue, very thick] {x^2/2 + 1};
\addplot[curveblue, very thick] {x^2/2 + 2};
\addplot[tangentred, thick, domain=0.55:1.45] {(x-1) + (0.5 - 1)};
\addplot[tangentred, thick, domain=0.55:1.45] {(x-1) + (0.5 + 0)};
\addplot[tangentred, thick, domain=0.55:1.45] {(x-1) + (0.5 + 1)};
\addplot[tangentred, thick, domain=0.55:1.45] {(x-1) + (0.5 + 2)};
\node[anchor=west, font=\small] at (axis cs:1.5,3.6) {same slope $=x$};
\node[anchor=south, font=\small, curveblue] at (axis cs:-0.7,3.1) {$F(x)+C$};
\end{axis}
\end{tikzpicture}
\caption{Antiderivatives of $f(x)=x$ are $\tfrac12 x^2+C$, a stack of vertical
shifts of one parabola. At any fixed $x$ all of them have the identical slope $x$
(short red tangents at $x=1$), which is why they share the same derivative.}
\label{ch21:fig-plusC}
\end{figure}

\begin{trapbox}
Two $+C$ pitfalls.

\emph{Forgetting it on an indefinite integral.} $\int 2x\,\dd x = x^2$ is
\textbf{incomplete}; the answer is $x^2+C$. You have named one antiderivative when
the question asked for all of them.

\emph{Writing it on a definite integral.} $\int_0^3 2x\,\dd x = 9$, a number --
no $C$. If you carry the constant through, it cancels in $F(3)-F(0)$ anyway, so
it is harmless but pointless. The $+C$ belongs to the indefinite world of
function-families, not the definite world of numbers.
\end{trapbox}

\begin{walkbox}
On your walk, picture the family of all antiderivatives as a deck of identical
curves stacked vertically. A definite integral asks you to read the \emph{rise}
between two $x$-values -- and rises are unaffected by which card of the deck you
read, because shifting the whole deck shifts both endpoints equally. Now ask
yourself: what \emph{one} extra piece of information would pin down a single card
out of the deck? (Answer: the value of $F$ at one point -- an ``initial
condition,'' the starting point for initial-value problems in differential equations.)
\end{walkbox}

\section{Worked evaluations}

\begin{recoverybox}
\textbf{Pause and think:} What happens if we integrate backward, with $b < a$?
Part~2 says $\int_a^b f(x)\,\dd x = F(b)-F(a)$. Notice this equals $-(F(a)-F(b)) = -\int_b^a f(x)\,\dd x$.
The evaluation recipe handles oriented limits automatically: running right to left flips the sign of the rise, mirroring the negative slice widths $\Delta x < 0$ in the underlying Riemann sum.
\end{recoverybox}

The procedure for evaluating a definite integral via Part 2 has two steps: (1) find any
antiderivative $F$ of the integrand; (2) compute $F(b)-F(a)$. The examples below
apply this directly.

\begin{example}\label{ch21:ex-poly}
$\displaystyle \int_1^3 (x^2 - 2x + 5)\,\dd x.$

\textbf{Antiderivative.} Reverse the power rule term by term:
$F(x)=\tfrac13 x^3 - x^2 + 5x$. (Check: $F'(x)=x^2-2x+5$. Good.)

\textbf{Evaluate.}
\[
  \eval{\tfrac13 x^3 - x^2 + 5x}{1}^{3}
  = \Big(9 - 9 + 15\Big) - \Big(\tfrac13 - 1 + 5\Big)
  = 15 - \tfrac{13}{3} = \frac{32}{3}.
\]
\end{example}

\begin{example}\label{ch21:ex-sin}
$\displaystyle \int_0^{\pi} \sin x\,\dd x.$

An antiderivative of $\sin x$ is $-\cos x$ (since $\ddx{-\cos x}=\sin x$). Thus
\[
  \eval{-\cos x}{0}^{\pi} = (-\cos\pi) - (-\cos 0) = \big(-(-1)\big) - \big(-(1)\big) = 1 - (-1) = 2.
\]
The area under one hump of the sine curve is exactly $2$ -- a clean number that
no amount of staring at Riemann rectangles would have suggested.
\end{example}

\begin{example}\label{ch21:ex-recip}
$\displaystyle \int_1^{e} \frac{1}{x}\,\dd x.$

The antiderivative of $1/x$ (for $x>0$) is $\ln x$. So
\[
  \eval{\ln x}{1}^{e} = \ln e - \ln 1 = 1 - 0 = 1.
\]
The area under the hyperbola $1/x$ from $1$ to $e$ is exactly $1$. This is, in
fact, the defining property of $e$: it is the upper limit that makes this area
equal to one.
\end{example}

\begin{example}\label{ch21:ex-exp}
$\displaystyle \int_0^{\ln 2} e^{x}\,\dd x.$

Since $e^x$ is its own antiderivative,
\[
  \eval{e^x}{0}^{\ln 2} = e^{\ln 2} - e^{0} = 2 - 1 = 1.
\]
\end{example}

\begin{example}\label{ch21:ex-sqrt}
$\displaystyle \int_0^{4} \sqrt{x}\,\dd x.$

Write $\sqrt{x}=x^{1/2}$. Reverse the power rule: raise the exponent by one to
$\tfrac32$ and divide, giving $F(x)=\tfrac{2}{3}x^{3/2}$. Then
\[
  \eval{\tfrac23 x^{3/2}}{0}^{4} = \tfrac23\cdot 4^{3/2} - 0
  = \tfrac23\cdot 8 = \frac{16}{3}.
\]
\end{example}

\begin{example}[Signed area, and why you sometimes split]\label{ch21:ex-signed}
$\displaystyle \int_0^{2\pi}\cos x\,\dd x$ versus the \emph{geometric} area
between $\cos x$ and the axis over the same interval.

The definite integral, using antiderivative $\sin x$:
\[
  \eval{\sin x}{0}^{2\pi} = \sin(2\pi)-\sin(0) = 0 .
\]
Zero -- the positive area of the first hump exactly cancels the negative area of
the trough (Figure~\ref{ch21:fig-signed}). If instead you want the \emph{total
geometric area} (everything counted positive), you must split at the sign
changes:
\[
  \int_0^{\pi/2}\!\cos x\,\dd x - \int_{\pi/2}^{3\pi/2}\!\cos x\,\dd x
  + \int_{3\pi/2}^{2\pi}\!\cos x\,\dd x
  = 1 + 2 + 1 = 4 .
\]

\begin{figure}[H]
\centering
\begin{tikzpicture}
\begin{axis}[
  width=12cm, height=5.6cm,
  axis lines=middle,
  xlabel={$x$}, ylabel={$\cos x$},
  xtick={1.5708,3.1416,4.7124,6.2832},
  xticklabels={$\frac{\pi}{2}$,$\pi$,$\frac{3\pi}{2}$,$2\pi$},
  ytick={-1,1}, xmin=-0.3, xmax=6.9, ymin=-1.5, ymax=1.5,
  domain=0:6.2832, samples=160, clip=true,
]
\addplot[draw=none, fill=areafill, opacity=0.6, domain=0:1.5708] {cos(deg(x))} \closedcycle;
\addplot[draw=none, fill=trapred, opacity=0.45, domain=1.5708:4.7124] {cos(deg(x))} \closedcycle;
\addplot[draw=none, fill=areafill, opacity=0.6, domain=4.7124:6.2832] {cos(deg(x))} \closedcycle;
\addplot[curveblue, very thick] {cos(deg(x))};
\node[font=\small] at (axis cs:0.78,0.45) {$+1$};
\node[font=\small] at (axis cs:3.14,-0.5) {$-2$};
\node[font=\small] at (axis cs:5.5,0.45) {$+1$};
\end{axis}
\end{tikzpicture}
\caption{Signed area of $\cos x$ over $[0,2\pi]$. The blue lobes above the axis
contribute $+1$ each and the red lobe below contributes $-2$, so the definite
integral is $1-2+1=0$. The total \emph{geometric} area counts every lobe
positively: $1+2+1=4$.}
\label{ch21:fig-signed}
\end{figure}
\end{example}

\begin{trapbox}
The single most common Part-2 disaster is integrating across a point where the
antiderivative \emph{blows up}. Consider $\int_{-1}^{1}\frac{1}{x^2}\,\dd x$.
Blindly using $F(x)=-\tfrac1x$ gives $-1-(1)=-2$, a \emph{negative} answer for
the area under a strictly positive function -- which contradicts the basic geometric
fact that the area under an everywhere-positive curve cannot be negative. The catch: $1/x^2$
is not continuous on $[-1,1]$; it explodes at $0$. Part 2 requires $f$ continuous
on the whole interval (and $F$ differentiable throughout). The integral
here is actually improper and diverges. \textbf{Always confirm the integrand is
continuous on $[a,b]$ before applying Part 2.}
\end{trapbox}

\begin{example}[Using Part 1 and Part 2 together]\label{ch21:ex-both}
Find the local maxima and minima of $\displaystyle A(x)=\int_0^x (t^2-4)\,\dd t$
without first computing $A$.

\textbf{Solution.} By Part 1, $A'(x)=x^2-4$. Critical points where $A'=0$ are
$x=\pm 2$. Since $A'(x)=x^2-4$ is negative on $(-2,2)$ and positive outside, $A$
decreases then increases: $x=2$ is a local minimum and $x=-2$ a local maximum.
(If you do want a value, Part 2 gives
$A(2)=\eval{\tfrac13 t^3-4t}{0}^{2}=\tfrac{8}{3}-8=-\tfrac{16}{3}$.) The lesson:
Part 1 turns questions about the \emph{shape} of an accumulation function into
ordinary derivative questions about the integrand.
\end{example}

\section{A short history of the bridge}

\begin{historybox}
Pieces of both halves were known long before anyone saw they fit together.
Archimedes (3rd century BCE, Syracuse) computed areas and volumes by his
``method of exhaustion'' (sandwiching a curved region between inscribed and
circumscribed polygons) -- essentially Riemann sums evaluated with remarkable
geometric skill, one shape at a time, with no general machinery. In the early
1600s, Pierre de Fermat, Bonaventura Cavalieri (using ``indivisibles,'' treating
a two-dimensional area as an infinite stack of one-dimensional lines), and John
Wallis found area formulas for $x^n$. On the other side, tangent-line methods
(the seeds of differentiation) were developed by Fermat, Ren\'e Descartes, and
especially Isaac Barrow, Newton's teacher and predecessor as Lucasian Professor
of Mathematics at Cambridge.

Barrow, in his 1670 \emph{Lectiones Geometricae} (specifically Lecture~X,
Proposition~11), came strikingly close to the general theorem. Using pure Euclidean
geometry without algebraic formulas, Barrow proved that if a curve is
constructed by accumulating the ordinates (the vertical heights $y$ of points on
a base curve), its tangent at any point is determined directly by the height of
the original curve. In modern terms, Barrow proved that differentiation inverts
accumulation. Yet Barrow did not develop this into a general computational framework.
Working entirely with static geometric diagrams without algebraic
equations or functional notation, he treated the proposition as merely one clever
geometric theorem among dozens, rather than a universal computational engine.

It took the algebraic and operational leap of \emph{Isaac Newton} (around 1666)
and \emph{Gottfried Wilhelm Leibniz} (in 1675) to recognize the Fundamental
Theorem as the unifying bridge of mathematics. Newton, motivated by physics and
celestial orbits, framed the calculus around ``fluxions'' (his term for the
instantaneous rates of change of flowing quantities) and ``fluents'' (the flowing
quantities themselves, such as accumulated position). Position is the accumulated
fluent of velocity. Leibniz, analyzing differential triangles and seeking a
universal symbolic language, engineered the supple $\dd x$ and elongated-S
$\int$ (for Latin \emph{summa}, sum) notation precisely so that $\dd$ and $\int$
would visibly undo each other on the page: $\int \dd F = F$ and
$\dd(\int f) = f\,\dd x$. Together, they transformed scattered geometric
triumphs into an algorithmic technology.

Two intense controversies followed. First was the bitter Newton--Leibniz priority
dispute: accusations of plagiarism, dueling royal academies, and decades of bad
blood that divided British and Continental mathematics for a century, though
modern historians confirm they discovered the bridge independently. Second was
the foundational crisis: both men computed with ``infinitesimals'' (quantities
treated as smaller than any positive real number yet strictly nonzero when
dividing to find slopes, only to be discarded as zero at the end). In 1734, the
philosopher Bishop George Berkeley published \emph{The Analyst}, criticizing this
shift between $\dd x \neq 0$ and $\dd x = 0$ as logically inconsistent. The objection
was fair: neither Newton nor Leibniz had a rigorous definition of an infinitesimal. It took Augustin-Louis Cauchy (1820s) with limits, Bernhard Riemann (1850s)
with partitioned sums, and Karl Weierstrass (1870s) with $\eps$--$\delta$
definitions to put the theorem on a rigorous foundation.

The theorem became the foundation for celestial mechanics, differential
equations, fluid dynamics, and physics. In geometry, it generalized
into Stokes' theorem on differential forms ($\int_{\partial \Omega}\omega = \int_\Omega \dd\omega$),
which unifies the Fundamental Theorem, Green's theorem, the divergence theorem,
and classical Stokes' theorem into a single line.
\end{historybox}

\section{The two parts, side by side}

\begin{focusbox}
A final consolidation. Memorize the \emph{shapes}, not just the symbols.

\textbf{Part 1 (the area-so-far has slope $f$):}
\[
  \frac{\dd}{\dd x}\int_a^x f(t)\,\dd t = f(x).
\]
\emph{Differentiating an accumulation returns the integrand.} This is the
existence half: it builds an antiderivative out of $f$, and it tells you the rate
at which area accumulates.

\textbf{Part 2 (evaluate by an antiderivative):}
\[
  \int_a^b f(x)\,\dd x = F(b)-F(a), \qquad F'=f.
\]
\emph{A definite integral is a net change in any antiderivative.} This is the
computational half: it converts the impossible-looking infinite sum into a
two-value subtraction.

They are two readings of one fact -- that accumulating and rate-measuring are
inverse processes -- and together they are the hinge on which all of integral
calculus turns. Every technique in the chapters ahead (substitution, parts,
partial fractions) is, at bottom, a trick for the \emph{first} step of Part 2:
\textbf{finding an antiderivative}. Once you have $F$, the theorem does the rest.
\end{focusbox}

\begin{applicationbox}
\textbf{Physics: position from velocity.} If $v(t)$ is velocity, then total
displacement over $[t_0,t_1]$ is $\int_{t_0}^{t_1} v(t)\,\dd t = s(t_1)-s(t_0)$,
because position $s$ is an antiderivative of $v$. The Fundamental Theorem is the
statement that ``displacement is the accumulation of velocity'' -- the area under
the speedometer trace is the distance traveled. Likewise impulse is the integral
of force, and work is the integral of force over distance. The theorem is why a
single antiderivative replaces an infinite tally of moment-by-moment
contributions.

\textbf{Probability and ML.} A probability density $p(x)$ has a cumulative
distribution $F(x)=\int_{-\infty}^x p(t)\,\dd t$; by Part 1, $F'=p$ -- the density
is the rate at which probability accumulates. In machine learning, the same
identity underlies how a model's loss accumulates from per-example contributions,
and why ``integrate the rate of change back'' recovers a net change.
\end{applicationbox}

\section{Reach of the Fundamental Theorem and the Liouville Barrier}

\begin{connectionbox}
Before the Fundamental Theorem, computing the area beneath a curve or the volume of a solid
required individual summation formulas for each specific curve, partitioning the domain and evaluating limits
of algebraic sums directly. The Fundamental Theorem unified these separate calculations into a clear principle:
\emph{to integrate, antidifferentiate}. Definite integrals across physics (work, impulse,
center of mass), geometry (arc length, surface area), and probability (cumulative distributions)
could now be evaluated directly from antiderivatives at boundary points, $F(b) - F(a)$.

Yet differentiation and antidifferentiation possess an asymmetric relationship: while
differentiation is an algorithmic mechanical process that always succeeds on elementary formulas,
antidifferentiation runs headlong into the celebrated \emph{Liouville barrier}. FTC Part~2 requires finding an
\emph{elementary antiderivative} $F$ such that $F' = f$. In 1835, Joseph Liouville proved
that simple, smooth, elementary functions exist whose antiderivatives \emph{cannot be expressed in terms of any finite
combination of elementary functions} (powers, exponentials, logarithms, trigonometric functions, and roots).
Canonical examples include:
\[
\begin{aligned}
  &\int e^{-x^2}\,\dd x \quad (\text{the Gaussian bell curve}), \qquad
  &&\int \frac{\sin x}{x}\,\dd x \quad (\text{the sine integral}), \\[4pt]
  &\int \frac{e^x}{x}\,\dd x \quad (\text{the exponential integral}), \qquad
  &&\int \sqrt{1 - k^2\sin^2 x}\,\dd x \quad (\text{the elliptic integral}).
\end{aligned}
\]
For all of these functions, FTC Part~1 guarantees that the area accumulation function
$A(x) = \int_0^x f(t)\,\dd t$ exists, is smooth, and satisfies $A'(x) = f(x)$. But FTC Part~2 is
powerless to evaluate $\int_a^b f(x)\,\dd x$ in closed elementary form, because no algebraic recipe
for $A(x)$ can ever be written down.

This tension sets the roadmap for the entire remainder of our study of integration:
\begin{itemize}
  \item \emph{Expanding the symbolic toolkit (Chapters 22--24):} When an elementary antiderivative
  exists, how do we systematically find it? Chapter~22 develops \emph{Integration by
  Substitution} (running the chain rule in reverse to untangle composite coordinates). Chapter~23
  introduces \emph{Integration by Parts} (inverting the product rule to disentangle products of
  dissimilar families). Chapter~24 builds \emph{Partial Fractions and Trigonometric Integrals} to
  completely solve all rational functions (ratios of polynomials).
  \item \emph{When formulas fail: Numerical Quadrature (Chapter 27):} When an integrand runs into
  Liouville's barrier---or when we must integrate discrete empirical measurements from sensors where
  no formula exists at all---we return to Riemann's original insight armed with modern numerical
  methods: the Trapezoidal Rule and Simpson's parabolic quadrature (numerical rules that approximate
  areas to high precision without ever needing an algebraic antiderivative).
\end{itemize}
\end{connectionbox}

In the next chapter,
we run the chain rule in reverse
to obtain \emph{Integration by Substitution}.

\section{Exercises}

\begin{exercise}[Mental checks]
Answer each question in your head or in a single sentence before calculating:
\begin{enumerate}[label=(\alph*)]
  \item If $f(t) \le 0$ everywhere on $[a,b]$, is the accumulation function $A(x)=\int_a^x f(t)\,\dd t$ increasing or decreasing?
  \item What is $\dv{}{x}\int_3^x \frac{\sin(\sqrt{1+t^5})}{2+\cos(t)}\,\dd t$ without doing any integration?
  \item Why does $\int_{-1}^1 \frac{1}{x^2}\,\dd x$ fail to equal $[\,-\tfrac{1}{x}\,]_{-1}^1 = -1 - (1) = -2$?
  \item If two students pick different antiderivatives of the same integrand, say $F_1(x) = x^2 + 3$ and $F_2(x) = x^2 - 17$, why do both obtain the exact same numerical answer for $\int_1^4 2x\,\dd x$?
\end{enumerate}
\end{exercise}

\begin{exercise}
Evaluate each definite integral with Part 2.
\begin{enumerate}[label=(\alph*)]
  \item $\displaystyle\int_0^2 (3x^2 - 4x + 1)\,\dd x$
  \item $\displaystyle\int_1^4 \frac{1}{\sqrt{x}}\,\dd x$
  \item $\displaystyle\int_0^{\pi/4}\sec^2 x\,\dd x$
  \item $\displaystyle\int_{-1}^{1} e^{2x}\,\dd x$
\end{enumerate}
\end{exercise}

\begin{exercise}
Differentiate each accumulation function using Part 1 (and the chain rule where
needed).
\begin{enumerate}[label=(\alph*)]
  \item $\displaystyle A(x)=\int_3^x \sqrt{1+t^4}\,\dd t$
  \item $\displaystyle B(x)=\int_0^{x^2} \cos(t)\,\dd t$
  \item $\displaystyle C(x)=\int_x^{5} \frac{\dd t}{1+t^2}$ (mind the sign)
  \item $\displaystyle D(x)=\int_{\sin x}^{\cos x} e^{t^2}\,\dd t$
\end{enumerate}
\end{exercise}

\begin{exercise}
Find the total \emph{geometric} area enclosed between the curve $y=x^2-1$ and the
$x$-axis on $[0,2]$. (Hint: the curve crosses the axis inside the interval, so
split and count each piece positively. Compare with the plain definite integral
$\int_0^2(x^2-1)\,\dd x$ and explain the difference.)
\end{exercise}

\begin{exercise}
Let $G(x)=\int_0^x (t^2 - 3t + 2)\,\dd t$. Using Part 1 only (do not first compute
$G$), find all critical points of $G$ and classify each as a local max or min.
Then, using Part 2, find the value of $G$ at each critical point.
\end{exercise}

\begin{exercise}[Conceptual]
A student writes $\displaystyle\int_{-2}^{2}\frac{1}{x^4}\,\dd x =
\eval{-\tfrac{1}{3x^3}}{-2}^{2} = -\tfrac{1}{24}-\tfrac{1}{24}=-\tfrac{1}{12}$.
Two things are wrong with this answer at once. Identify both: why is the numerical
sign absurd, and which hypothesis of Part 2 fails? What is the correct status of
this integral?
\end{exercise}

\begin{exercise}[Conceptual]
Explain, in your own words and without symbols where possible, why the constant
$C$ in $\int f\,\dd x = F+C$ cannot be detected by any definite integral
$\int_a^b f$. Connect your explanation to the picture of antiderivatives as
vertical shifts of one another.
\end{exercise}

\begin{exercise}
The error function is $\operatorname{erf}(x)=\frac{2}{\sqrt{\pi}}\int_0^x
e^{-t^2}\,\dd t$. Although $\operatorname{erf}$ has no elementary formula, compute
$\dv{}{x}\operatorname{erf}(x)$ and $\dv{}{x}\operatorname{erf}(x^2)$ exactly.
Explain why Part 1 lets you do this even though you cannot write down
$\operatorname{erf}$ in closed form.
\end{exercise}

\begin{exercise}[Mean Value Theorem for Integrals]
Suppose $f$ is continuous on $[a,b]$. Prove there exists $c\in[a,b]$ with
$\int_a^b f(x)\,\dd x = f(c)\,(b-a)$. (Hint: apply Part 1 to get the area function
$A$, then apply the ordinary Mean Value Theorem to $A$ on $[a,b]$. Interpret the
result geometrically: a continuous function actually attains its average value
somewhere.)
\end{exercise}

\begin{exercise}[Proof]
Prove that if $f$ is continuous on $[a,b]$ and $\int_a^x f(t)\,\dd t = 0$ for
\emph{every} $x\in[a,b]$, then $f$ is identically zero on $[a,b]$. (Hint:
differentiate the hypothesis using Part 1.) What does this say about the only
continuous function whose every accumulation vanishes?
\end{exercise}

\begin{exercise}[Synthesis]
A particle moves with velocity $v(t)=t^2-4t+3$ (meters per second) for
$t\in[0,4]$.
\begin{enumerate}[label=(\alph*)]
  \item Compute the net displacement $\int_0^4 v(t)\,\dd t$.
  \item Compute the total distance traveled (integrate $\abs{v}$, splitting at the
  times when $v=0$).
  \item Explain in one sentence why the two answers differ, in terms of the
  particle reversing direction. Which computation used signed area, and which used
  geometric area?
\end{enumerate}
\end{exercise}
